% GENERATED FILE. Edit canonical lessons, then run npm run generate:books.
% canonical-source-hash: fnv1a64:1d0510c903d56787
% canonical-lessons: GE-C530-kaemmen, GE-C530-gurgeln, GE-C530-schminken, GE-C530-bademantel, GE-C530-schlafanzug

\chapter{Morgens im Bad}
\label{ch:ge-morgens-im-bad}
\hlchaptermodality{530}

\begin{chapteropening}
\textbf{By the end of this chapter:} \emph{I can describe my routine in the bathroom.}

Complete the last lesson of chapter 530: I can describe my routine in the bathroom.
\end{chapteropening}

\section[kämmen]{kämmen — to comb}

\label{lesson:GE-C530-kaemmen}



\begin{quote}
Before the new one: say the German for tile, then the German for attic.
\end{quote}

\subsection*{You'll want to know: kämmen}
\textbf{kämmen} — \textquotedblleft{}to comb\textquotedblright{}.

\begin{quote}
Sie kämmt ihrer kleinen Tochter die Haare.
\end{quote}

\begin{quote}
\emph{She combs her little daughter's hair.}
\end{quote}

\emph{Sich kämmen} — to comb your hair. The comb is \emph{der Kamm}.

\begin{cousinweb}
\emph{kämmen} is inherited from Middle High German \emph{kemben} (EtymDB), built on the word that is today \emph{Kamm} \textquotedblleft{}comb\textquotedblright{}, from Old High German \emph{chamb} and Proto-Germanic \emph{kambaz}.
\end{cousinweb}

\begin{grammarlens}[title={What you've built}]
One more word for telling, explaining and describing at B1.
\end{grammarlens}

\subsection*{Your turn}
\begin{itemize}
\raggedright
  \item \emph{Say it:} \emph{kämmen}
  \item \emph{Say it:} \emph{kämmen}, once more
  \item \emph{Say it:} a sentence of your own with \emph{kämmen}
  \item \emph{Recall:} say the German for tile, then the German for attic, then say \emph{kämmen} again
\end{itemize}

\begin{tcolorbox}[breakable,colback=teal!4,colframe=teal!35!black,title={Before you move on}]
What does kämmen mean? (\textquotedblleft{}To comb\textquotedblright{}.) Say it once more.
\end{tcolorbox}
\section[gurgeln]{gurgeln — to gargle}

\label{lesson:GE-C530-gurgeln}



\begin{quote}
Before the new one: say the German for attic, then the German for to comb.
\end{quote}

\subsection*{You'll want to know: gurgeln}
\textbf{gurgeln} — \textquotedblleft{}to gargle\textquotedblright{}.

\begin{quote}
Bei Halsschmerzen gurgelt sie mit Salzwasser.
\end{quote}

\begin{quote}
\emph{When she has a sore throat, she gargles with salt water.}
\end{quote}

Regular: \emph{hat gegurgelt}. Also of water running away: \emph{Das Wasser gurgelt im Abfluss.}

\begin{cousinweb}
\emph{gurgeln} is built on \emph{Gurgel} \textquotedblleft{}throat\textquotedblright{}, which comes from Latin \emph{gurguliō} \textquotedblleft{}throat, windpipe\textquotedblright{} (EtymDB).
\end{cousinweb}

\begin{grammarlens}[title={What you've built}]
One more word for telling, explaining and describing at B1.
\end{grammarlens}

\subsection*{Your turn}
\begin{itemize}
\raggedright
  \item \emph{Say it:} \emph{gurgeln}
  \item \emph{Say it:} \emph{gurgeln}, once more
  \item \emph{Say it:} a sentence of your own with \emph{gurgeln}
  \item \emph{Recall:} say the German for attic, then the German for to comb, then say \emph{gurgeln} again
\end{itemize}

\begin{tcolorbox}[breakable,colback=teal!4,colframe=teal!35!black,title={Before you move on}]
What does gurgeln mean? (\textquotedblleft{}To gargle\textquotedblright{}.) Say it once more.
\end{tcolorbox}
\section[schminken]{schminken — to put on make-up (sich schminken)}

\label{lesson:GE-C530-schminken}



\begin{quote}
Before the new one: say the German for to comb, then the German for to gargle.
\end{quote}

\subsection*{You'll want to know: schminken}
\textbf{schminken} — \textquotedblleft{}to put on make-up (sich schminken)\textquotedblright{}.

\begin{quote}
Vor der Party schminkt sie sich im Bad.
\end{quote}

\begin{quote}
\emph{Before the party she does her make-up in the bathroom.}
\end{quote}

Mostly reflexive: \emph{sich schminken}. \emph{Die Schminke} — make-up.

\begin{cousinweb}
\emph{schminken} is built on \emph{Schminke} \textquotedblleft{}make-up\textquotedblright{}, which EtymDB links to Old High German \emph{smeichen} \textquotedblleft{}to smooth\textquotedblright{}; the details are not settled.
\end{cousinweb}

\begin{grammarlens}[title={What you've built}]
One more word for telling, explaining and describing at B1.
\end{grammarlens}

\subsection*{Your turn}
\begin{itemize}
\raggedright
  \item \emph{Say it:} \emph{schminken}
  \item \emph{Say it:} \emph{schminken}, once more
  \item \emph{Say it:} a sentence of your own with \emph{schminken}
  \item \emph{Recall:} say the German for to comb, then the German for to gargle, then say \emph{schminken} again
\end{itemize}

\begin{tcolorbox}[breakable,colback=teal!4,colframe=teal!35!black,title={Before you move on}]
What does schminken mean? (\textquotedblleft{}To put on make-up (sich schminken)\textquotedblright{}.) Say it once more.
\end{tcolorbox}
\section[der Bademantel]{der Bademantel — bathrobe, dressing gown}

\label{lesson:GE-C530-bademantel}



\begin{quote}
Before the new one: say the German for to gargle, then the German for to put on make-up.
\end{quote}

\subsection*{You'll want to know: der Bademantel}
\textbf{der Bademantel} — \textquotedblleft{}bathrobe, dressing gown\textquotedblright{}.

\begin{quote}
Am Sonntag frühstückt er im Bademantel.
\end{quote}

\begin{quote}
\emph{On Sundays he has breakfast in his bathrobe.}
\end{quote}

Plural: \emph{die Bademäntel}.

\begin{cousinweb}
\emph{Bademantel} is built on \emph{Bad} \textquotedblleft{}bath\textquotedblright{}, from Proto-Germanic \emph{baþą}, and \emph{Mantel} \textquotedblleft{}coat\textquotedblright{} (EtymDB lists both German words).
\end{cousinweb}

\begin{grammarlens}[title={What you've built}]
One more word for telling, explaining and describing at B1.
\end{grammarlens}

\subsection*{Your turn}
\begin{itemize}
\raggedright
  \item \emph{Say it:} \emph{der Bademantel}
  \item \emph{Say it:} \emph{der Bademantel}, once more
  \item \emph{Say it:} a sentence of your own with \emph{der Bademantel}
  \item \emph{Recall:} say the German for to gargle, then the German for to put on make-up, then say \emph{der Bademantel} again
\end{itemize}

\begin{tcolorbox}[breakable,colback=teal!4,colframe=teal!35!black,title={Before you move on}]
What does der Bademantel mean? (\textquotedblleft{}Bathrobe, dressing gown\textquotedblright{}.) Say it once more.
\end{tcolorbox}
\section[der Schlafanzug]{der Schlafanzug — pyjamas}

\label{lesson:GE-C530-schlafanzug}



\begin{quote}
Before the new one: say the German for to put on make-up, then the German for bathrobe.
\end{quote}

\subsection*{You'll want to know: der Schlafanzug}
\textbf{der Schlafanzug} — \textquotedblleft{}pyjamas\textquotedblright{}.

\begin{quote}
Die Kinder haben schon ihre Schlafanzüge an.
\end{quote}

\begin{quote}
\emph{The children already have their pyjamas on.}
\end{quote}

Plural: \emph{die Schlafanzüge}. One pair is \emph{ein Schlafanzug}.

\begin{cousinweb}
\emph{Schlafanzug} is built on \emph{schlafen} \textquotedblleft{}to sleep\textquotedblright{} and \emph{Anzug} \textquotedblleft{}suit\textquotedblright{} (EtymDB lists it), from \emph{anziehen} \textquotedblleft{}to put on\textquotedblright{}, built on \emph{ziehen} \textquotedblleft{}to pull\textquotedblright{}.
\end{cousinweb}

\begin{grammarlens}[title={What you've built}]
One more word for telling, explaining and describing at B1.
\end{grammarlens}

\subsection*{Your turn}
\begin{itemize}
\raggedright
  \item \emph{Say it:} \emph{der Schlafanzug}
  \item \emph{Say it:} \emph{der Schlafanzug}, once more
  \item \emph{Say it:} a sentence of your own with \emph{der Schlafanzug}
  \item \emph{Recall:} say the German for to put on make-up, then the German for bathrobe, then say \emph{der Schlafanzug} again
\end{itemize}

\begin{tcolorbox}[breakable,colback=teal!4,colframe=teal!35!black,title={Before you move on}]
What does der Schlafanzug mean? (\textquotedblleft{}Pyjamas\textquotedblright{}.) Say it once more.
\end{tcolorbox}
